\import{topology/topological-ring-f-adic.tex}
\import{topology/topological-ring-nilpotence.tex}

\section{Pseudouniformizers and Tate rings}

\begin{definition}\label{pseudouniformizer}
    $\pi$ is a pseudouniformizer of $R$ iff
        $R$ is a topological ring and $\pi\in\carrier[R]$ and
        $\pi$ is topologically nilpotent in $R$ and
        $\pi$ is a unit of $\MultiplicativeMonoid{R}$.
\end{definition}

\begin{proposition}\label{pseudouniformizer_ambient}
    Suppose $\pi$ is a pseudouniformizer of $R$.
    Then $R$ is a topological ring and $\pi\in\carrier[R]$.
\end{proposition}
\begin{proof}
    Follows by \obv{pseudouniformizer}.
\end{proof}

\begin{proposition}\label{pseudouniformizer_is_topologically_nilpotent}
    Suppose $\pi$ is a pseudouniformizer of $R$.
    Then $\pi$ is topologically nilpotent in $R$.
\end{proposition}
\begin{proof}
    Follows by \obv{pseudouniformizer}.
\end{proof}

\begin{proposition}\label{pseudouniformizer_is_unit}
    Suppose $\pi$ is a pseudouniformizer of $R$.
    Then $\pi$ is a unit of $\MultiplicativeMonoid{R}$.
\end{proposition}
\begin{proof}
    Follows by \obv{pseudouniformizer}.
\end{proof}

\begin{proposition}\label{pseudouniformizer_power_is_unit}
    Suppose $\pi$ is a pseudouniformizer of $R$.
    Let $n$ be a natural number.
    Then $\ringPower{R}{\pi}{n}$ is a unit of
        $\MultiplicativeMonoid{R}$.
\end{proposition}
\begin{proof}
    Hence $R$ is a ring by \obv{topological_ring_is_ring}.
    $\pi$ is a unit of $\MultiplicativeMonoid{R}$ by
        \obv{pseudouniformizer_is_unit}.
    Follows by \obv{ring_unit_power}.
\end{proof}

\begin{theorem}\label{pseudouniformizer_power_bounded}
    Let $R$ be a commutative topological ring.
    Suppose $\pi$ is a pseudouniformizer of $R$.
    Then $\pi$ is power bounded in $R$.
\end{theorem}
\begin{proof}
    $\pi$ is topologically nilpotent in $R$ by
        \obv{pseudouniformizer_is_topologically_nilpotent}.
    Follows by \obv{topologically_nilpotent_is_power_bounded}.
\end{proof}

\begin{corollary}\label{pseudouniformizer_in_power_bounded_carrier}
    Let $R$ be a commutative topological ring.
    Suppose $\pi$ is a pseudouniformizer of $R$.
    Then $\pi\in\PowerBoundedCarrier{R}$.
\end{corollary}
\begin{proof}
    $\pi\in\carrier[R]$ by \obv{pseudouniformizer_ambient}.
    $\pi$ is power bounded in $R$ by
        \obv{pseudouniformizer_power_bounded}.
    Follows by \obv{power_bounded_carrier_iff}.
\end{proof}

\begin{definition}\label{tate_ring}
    $R$ is a Tate ring iff
        $R$ is an f-adic ring and
        there exists $\pi$ such that $\pi$ is a pseudouniformizer of $R$.
\end{definition}

\begin{abbreviation}\label{commutative_tate_ring}
    $R$ is a commutative Tate ring iff $R$ is a Tate ring.
\end{abbreviation}

\begin{proposition}\label{tate_ring_is_f_adic_ring}
    Let $R$ be a Tate ring.
    Then $R$ is an f-adic ring.
\end{proposition}
\begin{proof}
    Follows by \obv{tate_ring}.
\end{proof}

\begin{classification}\label{tate_ring_is_commutative_topological_ring}
    Let $R$ be a Tate ring.
    Then $R$ is a commutative topological ring.
\end{classification}
\begin{proof}
    $R$ is an f-adic ring by \obv{tate_ring_is_f_adic_ring}.
    Follows by \obv{f_adic_ring_is_commutative_topological_ring}.
\end{proof}

\begin{proposition}\label{tate_ring_has_pseudouniformizer}
    Let $R$ be a Tate ring.
    Then there exists $\pi$ such that
        $\pi$ is a pseudouniformizer of $R$.
\end{proposition}
\begin{proof}
    Follows by \obv{tate_ring}.
\end{proof}

\begin{classification}\label{tate_ring_is_nonarchimedean}
    Let $R$ be a Tate ring.
    Then $R$ is a nonarchimedean topological ring.
\end{classification}
\begin{proof}
    $R$ is an f-adic ring by \obv{tate_ring_is_f_adic_ring}.
    Follows by \obv{f_adic_ring_is_nonarchimedean}.
\end{proof}

\begin{theorem}\label{tate_ring_has_power_bounded_pseudouniformizer}
    Let $R$ be a Tate ring.
    Then there exists $\pi\in\PowerBoundedCarrier{R}$ such that
        $\pi$ is a pseudouniformizer of $R$.
\end{theorem}
\begin{proof}
    $R$ is a commutative topological ring by
        \obv{tate_ring_is_commutative_topological_ring}.
    Take $\pi$ such that $\pi$ is a pseudouniformizer of $R$ by
        \obv{tate_ring_has_pseudouniformizer}.
    $\pi\in\PowerBoundedCarrier{R}$ by
        \obv{pseudouniformizer_in_power_bounded_carrier}.
\end{proof}

\begin{proposition}\label{unit_scalar_multiple_subring_power_neighbourhood}
    Let $R$ be a commutative topological ring and $A$ be a subring of $R$.
    Suppose $A$ is open in $R$.
    Suppose $a$ is a unit of $\MultiplicativeMonoid{R}$.
    Let $n$ be a natural number.
    Then $\fAdicPower{R}{
            \pointwiseProduct{\MultiplicativeMonoid{R}}{\{a\}}{A}}{
            \suc{n}}
        \in\neighbourhoods{\ringZero[R]}{R}$.
\end{proposition}
\begin{proof}
    $R$ is a topological ring by
        \obv{commutative_topological_ring_is_topological_ring}.
    $R$ is a commutative ring by
        \obv{commutative_topological_ring_is_commutative_ring}.
    Hence $R$ is a ring by \obv{commutative_ring_is_ring}.
    $a\in\carrier[\MultiplicativeMonoid{R}]$ by
        \obv{monoid_unit_in_carrier}.
    Hence $a\in\carrier[R]$ by
        \auto{multiplicative_monoid_fields}.
    $\suc{n}$ is a natural number by \obv{omega_successor}.
    $\suc{n}\neq\emptyset$ by \obv{suc_neq_emptyset}.
    $\additiveSubsetPower{R}{
            \pointwiseProduct{\MultiplicativeMonoid{R}}{\{a\}}{A}}{
            \suc{n}}
        =\pointwiseProduct{\MultiplicativeMonoid{R}}{
            \{\ringPower{R}{a}{\suc{n}}\}}{A}$ by
        \obv{additive_subset_power_of_scalar_multiple}.
    Hence $\fAdicPower{R}{
            \pointwiseProduct{\MultiplicativeMonoid{R}}{\{a\}}{A}}{
            \suc{n}}
        =\pointwiseProduct{\MultiplicativeMonoid{R}}{
            \{\ringPower{R}{a}{\suc{n}}\}}{A}$ by
        \obv{f_adic_power}.
    $\ringPower{R}{a}{\suc{n}}$ is a unit of
        $\MultiplicativeMonoid{R}$ by \obv{ring_unit_power}.
    $\pointwiseProduct{\MultiplicativeMonoid{R}}{
            \{\ringPower{R}{a}{\suc{n}}\}}{A}
        \in\neighbourhoods{\ringZero[R]}{R}$ by
        \obv{topological_ring_unit_left_multiple_subring_neighbourhood}.
\end{proof}

\begin{theorem}\label{topologically_nilpotent_scalar_multiple_powers_refine}
    Let $R$ be a commutative topological ring.
    Suppose $\pi$ is topologically nilpotent in $R$.
    Let $A$ be a subring of $R$ such that $\Bounded{R}{A}$.
    Let $N$ be a natural number such that $N\neq\emptyset$.
    Then for every $V\in\neighbourhoods{\ringZero[R]}{R}$ there exists
        $k\in\naturals$ such that
        $\fAdicPower{R}{
            \pointwiseProduct{\MultiplicativeMonoid{R}}{
                \{\ringPower{R}{\pi}{N}\}}{A}}{\suc{k}}\subseteq V$.
\end{theorem}
\begin{proof}
    $R$ is a topological ring by
        \obv{commutative_topological_ring_is_topological_ring}.
    $R$ is a commutative ring by
        \obv{commutative_topological_ring_is_commutative_ring}.
    Let $M=\MultiplicativeMonoid{R}$.
    $M$ is a monoid by \obv{multiplicative_monoid_is_exact_monoid}.
    Hence $M$ is a magma by \obv{monoid_is_magma}.
    $\carrier[M]=\carrier[R]$ by
        \auto{multiplicative_monoid_fields}.
    $A\subseteq\carrier[R]$ by \obv{subring_subseteq_carrier}.
    Hence $A\subseteq\carrier[M]$.
    $\pi\in\carrier[R]$ by
        \obv{topologically_nilpotent_element_ambient}.
    $\ringPower{R}{\pi}{N}\in\carrier[R]$ by
        \obv{ring_power_type}.
    $A$ is left bounded in $R$ by
        \obv{bounded_iff_left_bounded}.

    Fix $V\in\neighbourhoods{\ringZero[R]}{R}$.
    Take $W\in\neighbourhoods{\ringZero[R]}{R}$ such that
        $\pointwiseProduct{\MultiplicativeMonoid{R}}{W}{A}
            \subseteq V$ by
        \obv{left_bounded_in_neighbourhood}.
    Hence $\pointwiseProduct{M}{W}{A}\subseteq V$.
    $W\subseteq\carrier[R]$ by
        \auto{topological_ring_zero_neighbourhoods_is_filter,filter}.
    Hence $W\subseteq\carrier[M]$.
    Take $k\in\naturals$ such that for every $m\in\naturals$ such that
        $k\subseteq m$ we have $\ringPower{R}{\pi}{m}\in W$ by
        \auto{topologically_nilpotent_iff_eventually_in_neighbourhoods}.
    $\suc{k}$ is a natural number by \obv{omega_successor}.
    $\suc{k}\neq\emptyset$ by \obv{suc_neq_emptyset}.
    $\unions{N}$ is a natural number and $N=\suc{\unions{N}}$ by
        \obv{natural_nonzero_successor_of_union}.
    $\naturalMul{\unions{N}}{\suc{k}}$ is a natural number by
        \obv{natural_multiplication_closed}.
    $\naturalMul{N}{\suc{k}}
        =\naturalMul{\unions{N}}{\suc{k}}+\suc{k}$ by
        \auto{natural_multiplication_successor_left}.
    $\naturalMul{\unions{N}}{\suc{k}}+\suc{k}
        =\suc{k}+\naturalMul{\unions{N}}{\suc{k}}$ by
        \obv{natural_addition_commutative}.
    $\suc{k}\subseteq
        \suc{k}+\naturalMul{\unions{N}}{\suc{k}}$ by
        \obv{natural_addition_left_subseteq_sum}.
    Hence $\suc{k}\subseteq\naturalMul{N}{\suc{k}}$.
    $k\subseteq\suc{k}$ by \obv{subseteq_self_suc_intro}.
    Hence $k\subseteq\naturalMul{N}{\suc{k}}$ by
        \obv{subseteq_transitive}.
    $\naturalMul{N}{\suc{k}}$ is a natural number by
        \obv{natural_multiplication_closed}.
    Hence $\ringPower{R}{\pi}{\naturalMul{N}{\suc{k}}}\in W$.
    $\ringPower{R}{\pi}{\naturalMul{N}{\suc{k}}}
        =\ringPower{R}{\ringPower{R}{\pi}{N}}{\suc{k}}$ by
        \obv{ring_power_multiplication}.
    Hence $\ringPower{R}{\ringPower{R}{\pi}{N}}{\suc{k}}\in W$.
    Hence $\{\ringPower{R}{\ringPower{R}{\pi}{N}}{\suc{k}}\}
        \subseteq W$ by \auto{subseteq,singleton_iff}.
    $\pointwiseProduct{M}{
            \{\ringPower{R}{\ringPower{R}{\pi}{N}}{\suc{k}}\}}{A}
        \subseteq\pointwiseProduct{M}{W}{A}$ by
        \obv{pointwise_product_monotone,subseteq_refl}.
    Hence $\pointwiseProduct{M}{
            \{\ringPower{R}{\ringPower{R}{\pi}{N}}{\suc{k}}\}}{A}
        \subseteq V$ by \obv{subseteq_transitive}.
    $\additiveSubsetPower{R}{
            \pointwiseProduct{\MultiplicativeMonoid{R}}{
                \{\ringPower{R}{\pi}{N}\}}{A}}{\suc{k}}
        =\pointwiseProduct{\MultiplicativeMonoid{R}}{
            \{\ringPower{R}{\ringPower{R}{\pi}{N}}{\suc{k}}\}}{A}$ by
        \obv{additive_subset_power_of_scalar_multiple}.
    Hence $\fAdicPower{R}{
            \pointwiseProduct{\MultiplicativeMonoid{R}}{
                \{\ringPower{R}{\pi}{N}\}}{A}}{\suc{k}}
        =\pointwiseProduct{M}{
            \{\ringPower{R}{\ringPower{R}{\pi}{N}}{\suc{k}}\}}{A}$ by
        \obv{f_adic_power}.
    Follows by \auto.
\end{proof}

\section{The open bounded subring criterion}

\begin{theorem}\label{f_adic_ring_iff_open_bounded_subring}
    Let $R$ be a commutative topological ring.
    Suppose $\pi$ is a pseudouniformizer of $R$.
    Then $R$ is an f-adic ring iff there exists $A$ such that
        $A$ is a subring of $R$ and $A$ is open in $R$ and
        $\Bounded{R}{A}$.
\end{theorem}
\begin{proof}
    Show that if $R$ is an f-adic ring, then there exists $A$ such that
        $A$ is a subring of $R$ and $A$ is open in $R$ and
        $\Bounded{R}{A}$.
    \begin{subproof}
        Assume $R$ is an f-adic ring.
        Follows by \obv{f_adic_ring_has_open_bounded_subring}.
    \end{subproof}

    Show that if there exists $A$ such that $A$ is a subring of $R$ and
        $A$ is open in $R$ and $\Bounded{R}{A}$, then
        $R$ is an f-adic ring.
    \begin{subproof}
        Assume that there exists $A$ such that $A$ is a subring of $R$ and
            $A$ is open in $R$ and $\Bounded{R}{A}$.
        Take $A$ such that $A$ is a subring of $R$ and
            $A$ is open in $R$ and $\Bounded{R}{A}$.
        $R$ is a topological ring by
            \obv{commutative_topological_ring_is_topological_ring}.
        $R$ is a commutative ring by
            \obv{commutative_topological_ring_is_commutative_ring}.
        Let $M=\MultiplicativeMonoid{R}$.
        $\carrier[M]=\carrier[R]$ and $\mul[M]=\ringMul[R]$ by
            \auto{multiplicative_monoid_fields}.
        $A\subseteq\carrier[R]$ by \obv{subring_subseteq_carrier}.
        Hence $A\subseteq\carrier[M]$.
        $\subgroup{A}{\AdditiveGroup{R}}$ by
            \obv{subring_additive_subgroup}.
        $\pi\in\carrier[R]$ by \obv{pseudouniformizer_ambient}.
        Hence $\pi\in\carrier[M]$.
        $\pi$ is topologically nilpotent in $R$ by
            \obv{pseudouniformizer_is_topologically_nilpotent}.
        $\pi$ is a unit of $M$ by
            \obv{pseudouniformizer_is_unit}.

        $A\in\OpenSubrings{R}$ by \obv{open_subrings_iff}.
        Hence $A\in\neighbourhoods{\ringZero[R]}{R}$ by
            \obv{open_subring_iff_neighbourhood_of_zero}.
        Take $K\in\naturals$ such that for every $n\in\naturals$ such
            that $K\subseteq n$ we have $\ringPower{R}{\pi}{n}\in A$ by
            \auto{topologically_nilpotent_iff_eventually_in_neighbourhoods}.
        Let $N=\suc{K}$.
        $\suc{K}\in\naturals$ and $\suc{K}\neq\emptyset$ by
            \obv{omega_successor,suc_neq_emptyset}.
        Hence $N\in\naturals$ and $N\neq\emptyset$.
        $K\subseteq N$ by \obv{subseteq_self_suc_intro}.
        Hence $\ringPower{R}{\pi}{N}\in A$.
        Let $a=\ringPower{R}{\pi}{N}$.
        Hence $a\in A$.
        Hence $a\in\carrier[R]$ and $a\in\carrier[M]$ by
            \obv{subseteq}.
        $a$ is a unit of $\MultiplicativeMonoid{R}$ by
            \auto{pseudouniformizer_power_is_unit}.
        Hence $a$ is a unit of $M$.

        Let $U=\pointwiseProduct{M}{\{a\}}{A}$ and $T=\{a\}$.
        $\subgroup{
            \pointwiseProduct{\MultiplicativeMonoid{R}}{\{a\}}{A}}{
            \AdditiveGroup{R}}$ by
            \obv{left_scalar_multiple_is_additive_subgroup}.
        Hence $\subgroup{U}{\AdditiveGroup{R}}$.
        Hence $U\subseteq\carrier[\AdditiveGroup{R}]$ by
            \obv{subgroup_subseteq}.
        Hence $U\subseteq\carrier[R]$ by
            \auto{additive_group_fields}.
        $T$ is finite by \obv{singleton_is_finite}.
        $\ringOne[R]\in A$ by \obv{subring_one}.
        $\ringMul[R](a,\ringOne[R])\in U$ by
            \auto{pointwise_product_iff,singleton_iff}.
        $\ringMul[R](a,\ringOne[R])=a$ by
            \obv{ring_mul_one_right}.
        Hence $a\in U$.
        Hence $T\subseteq U$ by \auto{subseteq,singleton_iff}.

        For every $n\in\naturals$ we have
            $\fAdicPower{R}{
                \pointwiseProduct{\MultiplicativeMonoid{R}}{\{a\}}{A}}{
                \suc{n}}\in\neighbourhoods{\ringZero[R]}{R}$ by
            \obv{unit_scalar_multiple_subring_power_neighbourhood}.
        Hence for every $n\in\naturals$ we have
            $\fAdicPower{R}{U}{\suc{n}}
                \in\neighbourhoods{\ringZero[R]}{R}$ by \auto.

        For every $V\in\neighbourhoods{\ringZero[R]}{R}$ there exists
            $n\in\naturals$ such that
            $\fAdicPower{R}{
                \pointwiseProduct{\MultiplicativeMonoid{R}}{
                    \{\ringPower{R}{\pi}{N}\}}{A}}{\suc{n}}\subseteq V$ by
            \obv{topologically_nilpotent_scalar_multiple_powers_refine}.
        Hence for every $V\in\neighbourhoods{\ringZero[R]}{R}$ there
            exists $n\in\naturals$ such that
                $\fAdicPower{R}{U}{\suc{n}}\subseteq V$ by \auto.

        Let $Z=\pointwiseProduct{M}{
            \{\ringMul[R](a,a)\}}{A}$.
        $\additiveSubsetProduct{R}{\{a\}}{
                \pointwiseProduct{\MultiplicativeMonoid{R}}{\{a\}}{A}}
            =\pointwiseProduct{\MultiplicativeMonoid{R}}{
                \{\ringMul[R](a,a)\}}{A}$ by
            \obv{additive_subset_product_singleton_scalar_multiple}.
        Hence $\additiveSubsetProduct{R}{T}{U}=Z$.
        $\additiveSubsetProduct{R}{
                \pointwiseProduct{\MultiplicativeMonoid{R}}{\{a\}}{A}}{
                \pointwiseProduct{\MultiplicativeMonoid{R}}{\{a\}}{A}}
            =\pointwiseProduct{\MultiplicativeMonoid{R}}{
                \{\ringMul[R](a,a)\}}{A}$ by
            \obv{additive_subset_product_of_scalar_multiples}.
        Hence $\additiveSubsetProduct{R}{U}{U}=Z$.
        Hence $\additiveSubsetProduct{R}{T}{U}
            =\additiveSubsetProduct{R}{U}{U}$.

        $\pointwiseProduct{\MultiplicativeMonoid{R}}{
                \{\ringMul[R](a,a)\}}{A}
            \subseteq\pointwiseProduct{\MultiplicativeMonoid{R}}{
                \{a\}}{A}$ by
            \obv{scalar_square_subring_subset_scalar_multiple}.
        Hence $Z\subseteq U$.
        Hence $\additiveSubsetProduct{R}{U}{U}\subseteq U$.
        $\fAdicPresentation{R}{U}{T}$ by
            \obv{f_adic_presentation}.
        Follows by \auto{f_adic_ring}.
    \end{subproof}
    Follows by \auto.
\end{proof}
